\chapter{Intégration sur un segment}\label{ch:b1:integration}

L'\omterm{thm:b1:integration:def}{intégrale} du volume de Terminale reposait sur des aires prises
intuitivement. Ce chapitre la construit~: d'abord pour les \omterm{def:b1:integration:step}{fonctions
en escalier}, où l'\omterm{thm:b1:integration:def}{intégrale} est une somme finie, puis pour les
fonctions \omterm{def:b1:continuity:continuous}{continues} (et \omterm{def:b1:continuity:continuous}{continues} par morceaux) par approximation
uniforme --- l'endroit où le \omterm{thm:b1:continuity:heine}{théorème de Heine}
(\cref{thm:b1:continuity:heine}) montre toute son utilité. Le
\omterm{thm:b1:integration:ftc}{théorème fondamental de l'analyse} relie ensuite la construction aux
primitives, et les \omterm{thm:b1:integration:riemann}{sommes de Riemann} la relient aux moyennes
discrètes.

Dans tout le chapitre, $a < b$ sont des réels.

\section{Fonctions en escalier}

\begin{definition}\label{def:b1:integration:step}
$\varphi \colon \intcc{a}{b} \to \R$ est une \emph{fonction en
escalier}\index{fonction en escalier} lorsqu'il existe une
subdivision $a = x_0 < x_1 < \dots < x_n = b$ telle que $\varphi$
soit constante, égale à $c_i$, sur chaque \omterm{prop:b1:reals:intervals}{intervalle} \omterm{def:b1:topology:open}{ouvert}
$\intoo{x_{i-1}}{x_i}$ (les valeurs aux nœuds sont quelconques). Son
\omterm{thm:b1:integration:def}{intégrale} est
\[
\int_a^b \varphi = \sum_{i=1}^{n} c_i\,(x_i - x_{i-1}),
\]
indépendante de la subdivision choisie (on raffine deux subdivisions
par leur subdivision commune~: chaque membre est inchangé par
raffinement).
\end{definition}

\begin{proposition}\label{prop:b1:integration:stepprops}
Sur les \omterm{def:b1:integration:step}{fonctions en escalier}, l'\omterm{thm:b1:integration:def}{intégrale} est linéaire, croissante
($\varphi \leq \psi \implies \int\varphi \leq \int\psi$), et vérifie
la relation de Chasles $\int_a^b = \int_a^c + \int_c^b$ pour $a < c <
b$.
\end{proposition}

\begin{proof}
Le moteur est l'\emph{invariance par raffinement}, énoncée dans la
définition~: insérer un nœud supplémentaire $t \in
\intoo{x_{i-1}}{x_i}$ dans une subdivision remplace le terme
$c_i(x_i - x_{i-1})$ par $c_i(t - x_{i-1}) + c_i(x_i - t)$ --- le
même nombre --- de sorte que l'\omterm{thm:b1:integration:def}{intégrale} est inchangée par tout
raffinement fini. Prenons maintenant $\varphi$ de subdivision
$\sigma$ et $\psi$ de subdivision $\sigma'$~: sur le raffinement
commun $\sigma \cup \sigma'$, toutes deux sont des \omterm{def:b1:integration:step}{fonctions en
escalier} ayant les \emph{mêmes} nœuds, et sur chaque morceau
$\varphi + \lambda\psi$ est constante égale à $c_i + \lambda
d_i$~: la linéarité se ramène à la linéarité des sommes finies.
Croissance~: $c_i \leq d_i$ sur chaque morceau donne $\sum c_i
\Delta_i \leq \sum d_i \Delta_i$ (les longueurs $\Delta_i \geq
0$). Chasles~: insérer le nœud $c$ et couper la somme en ce point.
\end{proof}

\section{Intégrale d'une fonction continue}

\begin{theorem}[Approximation uniforme]\label{thm:b1:integration:approx}
Soit $f$ \omterm{def:b1:continuity:continuous}{continue} sur $\intcc{a}{b}$. Pour tout $\varepsilon > 0$, il
existe des \omterm{def:b1:integration:step}{fonctions en escalier} $\varphi, \psi$ telles que
\[
\varphi \leq f \leq \psi
\qquad\text{et}\qquad
\psi - \varphi \leq \varepsilon \text{ sur } \intcc{a}{b}.
\]
\end{theorem}

\begin{proof}
D'après le \omterm{thm:b1:continuity:heine}{théorème de Heine} (\cref{thm:b1:continuity:heine}), $f$ est
\omterm{def:b1:continuity:uniform}{uniformément continue}~: choisissons $\delta$ associé à $\varepsilon$,
et une subdivision de pas $< \delta$ (régulière, disons, avec $n >
\frac{b - a}{\delta}$ morceaux). Sur chaque morceau \omterm{def:b1:topology:closed}{fermé}
$\intcc{x_{i-1}}{x_i}$, $f$ atteint un minimum $m_i$ et un maximum
$M_i$ (\cref{thm:b1:continuity:evt}), et $M_i - m_i \leq \varepsilon$
(les deux points extrémaux sont distants de moins de $\delta$).
Posons $\varphi = m_i$ et $\psi = M_i$ sur $\intoo{x_{i-1}}{x_i}$ (et
$\varphi = \psi = f$ aux nœuds).
\end{proof}

\begin{theorem}[Définition de l'intégrale]\label{thm:b1:integration:def}
Soit $f$ \omterm{def:b1:continuity:continuous}{continue} sur $\intcc{a}{b}$. Les deux nombres
\[
I_-(f) = \sup\Bigl\{\int_a^b \varphi : \varphi \text{ en escalier},\
\varphi \leq f\Bigr\},
\qquad
I_+(f) = \inf\Bigl\{\int_a^b \psi : \psi \text{ en escalier},\ \psi \geq
f\Bigr\}
\]
sont égaux~; leur valeur commune est l'\emph{intégrale}\index{intégrale}
$\int_a^b f$ (également notée $\int_a^b f(t)\,\dd t$). Elle coïncide
avec la notion précédente sur les \omterm{def:b1:integration:step}{fonctions en escalier}, et s'étend
aux fonctions \omterm{def:b1:continuity:continuous}{continues} par morceaux en découpant $\intcc{a}{b}$ aux
discontinuités (Chasles y servant de définition).
\end{theorem}

\begin{proof}
Les deux \omterm{def:b1:logic:sets}{ensembles} sont non vides ($f$ est bornée) et toute \omterm{thm:b1:integration:def}{intégrale}
inférieure en escalier est $\leq$ à toute \omterm{thm:b1:integration:def}{intégrale} supérieure
(croissance sur les \omterm{def:b1:integration:step}{fonctions en escalier})~: donc $I_-(f) \leq
I_+(f)$. D'après le \cref{thm:b1:integration:approx}, pour tout
$\varepsilon$ il existe un couple avec $\int\psi - \int\varphi \leq
\varepsilon(b - a)$~: la \omterm{def:b1:reals:bounds}{borne supérieure} et la \omterm{def:b1:reals:bounds}{borne inférieure} sont
enserrées l'une contre l'autre, $I_- = I_+$.
\end{proof}

\begin{example}[Continue par morceaux, sans drame]\label{ex:b1:integration:piecewise}
La \omterm{thm:b1:reals:floor}{fonction partie entière} sur $\intcc{0}{3}$ est une \omterm{def:b1:integration:step}{fonction en
escalier} déguisée~: en coupant à ses sauts,
\[
\int_0^3 \lfloor t \rfloor\,\dd t
= \int_0^1 0 + \int_1^2 1 + \int_2^3 2 = 0 + 1 + 2 = 3 ,
\]
et les valeurs \emph{aux} points de saut $1, 2$ n'ont aucune
importance~: modifier une fonction en un nombre fini de points ne
change aucune \omterm{thm:b1:integration:def}{intégrale} (les \omterm{def:b1:integration:step}{fonctions en escalier} qui l'encadrent
n'en sont pas affectées). C'est tout le contenu de l'extension
«~\omterm{def:b1:continuity:continuous}{continue} par morceaux~»~: couper aux discontinuités, en nombre
fini, intégrer chaque morceau \omterm{def:b1:continuity:continuous}{continu}, additionner --- Chasles
comme définition.
\end{example}

\begin{example}[La définition calcule, une fois]\label{ex:b1:integration:defcompute}
Soit $f(x) = x$ sur $\intcc{0}{1}$, découpé en $n$ morceaux
égaux. Les meilleures \omterm{def:b1:integration:step}{fonctions en escalier} constantes sur les
morceaux sont $\varphi = \frac{k-1}{n}$ et $\psi = \frac kn$ sur
le $k$-ième morceau, avec
\[
\int_0^1 \varphi = \sum_{k=1}^{n} \frac{k-1}{n}\cdot\frac1n
= \frac{n-1}{2n},
\qquad
\int_0^1 \psi = \sum_{k=1}^{n} \frac{k}{n}\cdot\frac1n
= \frac{n+1}{2n} .
\]
Toute \omterm{thm:b1:integration:def}{intégrale} inférieure est $\leq I_-(f) \leq I_+(f) \leq$ à
toute \omterm{thm:b1:integration:def}{intégrale} supérieure, donc $\frac{n-1}{2n} \leq I_-(f) \leq
I_+(f) \leq \frac{n+1}{2n}$ pour tout $n$~: les deux sont
enserrées sur $\frac12$, et $\int_0^1 x\,\dd x = \frac12$
directement à partir de la définition. L'idée à retenir~: c'est la
première et la dernière fois que nous intégrons à partir de la
définition --- le théorème fondamental ci-dessous remplace tous
ces calculs par la lecture d'une primitive, ce qui est tout
l'intérêt économique de ce chapitre.
\end{example}

\begin{theorem}[Propriétés]\label{thm:b1:integration:props}
Pour $f, g$ \omterm{def:b1:continuity:continuous}{continues} (ou \omterm{def:b1:continuity:continuous}{continues} par morceaux) sur $\intcc{a}{b}$
et $\lambda \in \R$~:
\begin{enumerate}
  \item linéarité~: $\int (f + \lambda g) = \int f + \lambda \int g$~;
  \item croissance~: $f \leq g \implies \int_a^b f \leq \int_a^b g$~;
        et $\bigl|\int_a^b f\bigr| \leq \int_a^b \abs f \leq (b -
        a)\, \sup\abs f$~;
  \item Chasles~: $\int_a^b = \int_a^c + \int_c^b$ (avec la
        convention $\int_b^a = -\int_a^b$, valable quel que soit
        l'ordre des bornes)~;
  \item positivité stricte~: si $f$ est \omterm{def:b1:continuity:continuous}{continue}, $f \geq 0$ et
        $\int_a^b f = 0$, alors $f = 0$ partout sur
        $\intcc{a}{b}$.
\end{enumerate}
\end{theorem}

\begin{proof}
(1)--(3) passent des \omterm{def:b1:integration:step}{fonctions en escalier} à la limite par la
définition en sup/inf. La linéarité mérite qu'on la détaille une
fois~: soit $\varepsilon > 0$~; encadrons $\varphi_f \leq f \leq
\psi_f$ et $\varphi_g \leq g \leq \psi_g$ avec des écarts $\leq
\varepsilon$ (\cref{thm:b1:integration:approx}). Pour $\lambda
\geq 0$, $\varphi_f + \lambda\varphi_g \leq f + \lambda g \leq
\psi_f + \lambda\psi_g$ est un encadrement par des \omterm{def:b1:integration:step}{fonctions en
escalier} d'écart $\leq (1 + \lambda)\varepsilon$, et ses
\omterm{thm:b1:integration:def}{intégrales} en escalier valent $\int \varphi_f +
\lambda\int\varphi_g$, etc.\
(\cref{prop:b1:integration:stepprops})~: en faisant $\varepsilon
\to 0$, on enserre $\int(f + \lambda g)$ sur $\int f +
\lambda\int g$. Pour $\lambda < 0$, la multiplication par
$\lambda$ \emph{renverse} l'encadrement de $g$ --- la \omterm{def:b1:integration:step}{fonction en
escalier} inférieure de $\lambda g$ est $\lambda \psi_g$ --- et le
même encadrement fonctionne, les rôles échangés. La majoration
$\abs{\int f} \leq \int \abs f$ vient de $-\abs f \leq f \leq
\abs f$ et de la croissance.

(4) Par contraposée~: si $f(x_0) = m > 0$, la \omterm{def:b1:continuity:continuous}{continuité} fournit
un \omterm{prop:b1:reals:intervals}{sous-intervalle} de longueur $\eta > 0$ sur lequel $f \geq
\frac m2$~; la \omterm{def:b1:integration:step}{fonction en escalier} valant $\frac m2$ là et $0$
ailleurs est $\leq f$, donc $\int f \geq \frac{m\eta}{2} > 0$.
\end{proof}

\begin{example}[Chasles à l'œuvre~: intégrales de valeurs absolues]\label{ex:b1:integration:chasleswork}
Pour intégrer une valeur absolue, on coupe là où le signe change.
\[
\int_0^2 \abs{x - 1}\,\dd x
= \int_0^1 (1 - x)\,\dd x + \int_1^2 (x - 1)\,\dd x
= \frac12 + \frac12 = 1 ,
\]
et, en coupant $\intcc{0}{2\pi}$ en $\pi$~:
\[
\int_0^{2\pi} \abs{\sin t}\,\dd t
= \int_0^{\pi} \sin t\,\dd t - \int_{\pi}^{2\pi} \sin t\,\dd t
= 2 + 2 = 4 ,
\]
alors que $\int_0^{2\pi} \sin t\,\dd t = 0$~: les compensations
sont bien réelles, et c'est pourquoi l'énoncé de positivité
stricte (\cref{thm:b1:integration:props}~(4)) porte l'hypothèse $f
\geq 0$ --- sans elle, une \omterm{thm:b1:integration:def}{intégrale} nulle ne prouve rien sur
$f$. L'idée à retenir~: $\int \abs f$ mesure une \emph{aire},
$\int f$ mesure un \emph{bilan signé}~; l'inégalité $\abs{\int f}
\leq \int\abs f$ est le relevé exact de ce que les compensations
peuvent détruire.
\end{example}

\section{Le théorème fondamental de l'analyse}

\begin{theorem}[Théorème fondamental de l'analyse]\label{thm:b1:integration:ftc}
Soit $f$ \omterm{def:b1:continuity:continuous}{continue} sur un \omterm{prop:b1:reals:intervals}{intervalle} $I$ et $a \in I$. La fonction
\[
F(x) = \int_a^x f(t)\, \dd t
\]
est de classe $C^1$ sur $I$, avec $F' = f$~: toute fonction \omterm{def:b1:continuity:continuous}{continue}
sur un \omterm{prop:b1:reals:intervals}{intervalle} admet des primitives. Par conséquent, pour
\emph{toute} primitive $G$ de $f$~:\index{théorème fondamental de l'analyse}
\[
\int_a^b f(t)\,\dd t = G(b) - G(a) .
\]
\end{theorem}

\begin{proof}
Fixons $x_0 \in I$ et $\varepsilon > 0$~; la \omterm{def:b1:continuity:continuous}{continuité} en $x_0$
fournit $\delta$ tel que $\abs{f(t) - f(x_0)} \leq \varepsilon$ dès
que $\abs{t - x_0} \leq \delta$. Pour $0 < \abs{h} \leq \delta$ (et
$x_0 + h \in I$), Chasles donne
\[
\frac{F(x_0 + h) - F(x_0)}{h} - f(x_0)
= \frac 1h \int_{x_0}^{x_0+h} \bigl(f(t) - f(x_0)\bigr)\dd t ,
\]
dont la valeur absolue est $\leq \frac{1}{\abs h}\cdot \abs h\,
\varepsilon = \varepsilon$ (majoration (2), valable quel que soit
l'ordre des bornes). Donc $F'(x_0) = f(x_0)$~; $F' = f$ est
\omterm{def:b1:continuity:continuous}{continue}~: $F$ est $C^1$. Si $G' = f$ également, alors $(G - F)' = 0$
sur l'\omterm{prop:b1:reals:intervals}{intervalle}, donc $G = F + c$
(\cref{cor:b1:derivative:monotone}), et $G(b) - G(a) = F(b) - F(a) =
\int_a^b f$.
\end{proof}

\begin{example}[La symétrie avant le calcul]\label{ex:b1:integration:parity}
Sur un \omterm{prop:b1:reals:intervals}{intervalle} symétrique, la parité fait le travail~: si $f$
est impaire, le \omterm{thm:b1:integration:parts}{changement de variable} $t \mapsto -t$ envoie
$\int_{-a}^{0} f$ sur $-\int_0^a f$, donc
\[
\int_{-a}^{a} f(t)\,\dd t = 0 ;
\qquad\text{si $f$ est paire,}\quad
\int_{-a}^{a} f = 2\int_0^a f .
\]
Ainsi $\int_{-1}^{1} \frac{t^3\cos t}{1 + t^4}\,\dd t = 0$ sans
qu'aucune primitive soit en vue (l'intégrande est impaire), et
\[
\int_{-\pi}^{\pi} t^2\cos t\,\dd t = 2\int_0^\pi t^2\cos t\,
\dd t .
\]
Vérifier les symétries avant de sortir les techniques~:
l'\omterm{thm:b1:integration:def}{intégrale} la plus rapide est celle qu'on ne calcule jamais.
\end{example}

\begin{example}[Reconnaître une dérivée au premier coup d'œil]\label{ex:b1:integration:halfangle}
Calculons $\displaystyle\int_0^{\pi/2} \frac{\dd x}{1 + \cos x}$.
L'identité de l'angle moitié $1 + \cos x = 2\cos^2\frac x2$
transforme l'intégrande en $\frac{1}{2}\bigl(1 + \tan^2\frac
x2\bigr)$, qui est exactement la \omterm{def:b1:derivative:def}{dérivée} de $\tan\frac x2$~:
\[
\int_0^{\pi/2} \frac{\dd x}{1 + \cos x}
= \Bigl[\tan\frac x2\Bigr]_0^{\pi/2} = \tan\frac\pi4 = 1 .
\]
Aucune machinerie de \omterm{thm:b1:integration:parts}{changement de variable} n'a été nécessaire ---
seulement le réflexe de lire une intégrande comme la \omterm{def:b1:derivative:def}{dérivée} de
quelqu'un, le théorème fondamental faisant le reste. (L'outil
systématique qui gouverne de telles \omterm{thm:b1:integration:def}{intégrales} trigonométriques,
le \omterm{thm:b1:integration:parts}{changement de variable} $t = \tan\frac x2$, appartient à la
boîte à outils standard bâtie sur
le \cref{thm:b1:integration:parts}~(2).)
\end{example}

\begin{example}[Fonctions définies par une intégrale]\label{ex:b1:integration:definedby}
Le théorème fondamental fabrique des fonctions. Posons
\[
F(x) = \int_0^x \eu^{-t^2}\,\dd t .
\]
Aucune combinaison de fonctions classiques n'a pour \omterm{def:b1:derivative:def}{dérivée}
$\eu^{-t^2}$ (un théorème de Liouville, admis)~; pourtant $F$
existe, est $C^1$ avec $F'(x) = \eu^{-x^2} > 0$, strictement
croissante, impaire (\omterm{thm:b1:integration:parts}{changement de variable} $t \mapsto -t$), et
bornée~: pour $x \geq 1$,
\[
F(x) - F(1) = \int_1^x \eu^{-t^2}\dd t
\leq \int_1^x \eu^{-t}\dd t \leq \eu^{-1} ,
\]
donc $F \leq F(1) + \eu^{-1} \leq 1 + \eu^{-1}$. (La limite exacte,
$\frac{\sqrt\pi}{2}$, se calcule par \omterm{thm:b1:integration:def}{intégrales} doubles dans le
volume de Licence 3.) Dérivation d'une composée pour des bornes
mobiles~: $\frac{\dd}{\dd x}\int_x^{x^2} \eu^{-t^2}\dd t =
2x\,\eu^{-x^4} - \eu^{-x^2}$. L'idée à retenir~: l'intégration
\emph{crée} de nouvelles fonctions à partir des anciennes, avec
toutes leurs propriétés lisibles sur l'intégrande --- la primitive
qu'on ne sait pas écrire reste une fonction que l'on maîtrise
complètement.
\end{example}

\begin{example}[Estimer sans calculer]\label{ex:b1:integration:estimate}
Les \omterm{thm:b1:integration:def}{intégrales} $R_n = \int_0^1 \frac{t^n}{1 + t}\,\dd t$ n'ont
aucune forme close agréable, et pourtant la croissance les cerne
précisément~: sur $\intcc{0}{1}$, $\frac12 \leq \frac{1}{1+t} \leq
1$, donc
\[
\frac{1}{2(n+1)} = \frac12\int_0^1 t^n\,\dd t
\;\leq\; R_n \;\leq\; \int_0^1 t^n \,\dd t = \frac{1}{n+1} :
\]
l'ordre exact de décroissance ($R_n \sim$ un multiple de
$\frac1n$, en fait $R_n \sim \frac{1}{2n}$) en deux lignes et sans
primitive. Les devoirs maison de ce chapitre et du suivant
reposent exactement sur de tels encadrements --- le premier
réflexe de l'analyste devant une \omterm{thm:b1:integration:def}{intégrale} doit être de la
\emph{majorer}, et seulement ensuite, s'il le faut, de la
calculer.
\end{example}

\begin{example}[Valeurs moyennes]\label{ex:b1:integration:average}
La \emph{moyenne} d'une fonction \omterm{def:b1:continuity:continuous}{continue} $f$ sur $\intcc{a}{b}$
est $\frac{1}{b-a}\int_a^b f$. Pour l'arche du sinus~:
\[
\frac{1}{\pi}\int_0^\pi \sin t\,\dd t
= \frac{1}{\pi}\bigl[-\cos t\bigr]_0^\pi = \frac{2}{\pi}
\approx 0.637 :
\]
une arche positive complète a pour moyenne non pas $\frac12$ mais
$\frac2\pi$ --- la courbe passe plus de temps en hauteur qu'un
triangle ne le ferait. D'après l'\cref{exo:b1:integration:11}
(formule de la moyenne, $g = 1$), la moyenne est une
\emph{valeur}~: $\sin c = \frac2\pi$ pour un certain $c \in
\intoo{0}{\pi}$. Et d'après les \omterm{thm:b1:integration:riemann}{sommes de Riemann} de ce chapitre,
la moyenne est la limite des moyennes ordinaires de $n$
échantillons --- le pont entre la moyenne discrète de données et
la moyenne \omterm{def:b1:continuity:continuous}{continue} d'un signal, et c'est ainsi que l'\omterm{thm:b1:integration:def}{intégrale}
entre en physique.
\end{example}

\begin{theorem}[Intégration par parties~; changement de variable]\label{thm:b1:integration:parts}
\begin{enumerate}
  \item Si $u, v$ sont $C^1$ sur $\intcc{a}{b}$~:\index{intégration par
        parties}
        \[
        \int_a^b u'v = \bigl[uv\bigr]_a^b - \int_a^b uv' .
        \]
  \item Si $\varphi$ est $C^1$ sur $\intcc{\alpha}{\beta}$ et $f$
        \omterm{def:b1:continuity:continuous}{continue} sur $\varphi(\intcc{\alpha}{\beta})$~:
        \index{changement de variable}
        \[
        \int_{\alpha}^{\beta} f\bigl(\varphi(t)\bigr)\,\varphi'(t)\,
        \dd t = \int_{\varphi(\alpha)}^{\varphi(\beta)} f(x)\, \dd x .
        \]
\end{enumerate}
\end{theorem}

\begin{proof}
(1) $(uv)' = u'v + uv'$~; on intègre sur $\intcc{a}{b}$ et on applique
le théorème fondamental à la fonction $C^1$ $uv$.

(2) Soit $F$ une primitive de $f$ sur l'\omterm{prop:b1:reals:intervals}{intervalle} image
(\cref{thm:b1:integration:ftc}). Alors $(F \circ \varphi)' = (f
\circ \varphi)\,\varphi'$ (\omterm{def:b1:derivative:def}{dérivée} d'une composée), donc les deux
membres valent $F(\varphi(\beta)) - F(\varphi(\alpha))$.
\end{proof}

\begin{example}\label{ex:b1:integration:parts}
$\int_0^1 t\,\eu^t \dd t = \bigl[t\,\eu^t\bigr]_0^1 - \int_0^1
\eu^t\dd t = \eu - (\eu - 1) = 1$. Et avec le \omterm{thm:b1:integration:parts}{changement de variable}
$x = \sin t$ ($t \in \intcc{0}{\frac\pi2}$)~:
\[
\int_0^1 \sqrt{1 - x^2}\, \dd x
= \int_0^{\pi/2} \cos t \cdot \cos t \, \dd t
= \int_0^{\pi/2} \frac{1 + \cos 2t}{2}\, \dd t = \frac\pi4 ,
\]
--- un quart du disque unité, comme la géométrie l'exige. (La
linéarisation de la \cref{met:b1:complex:trig} à l'œuvre.)
\end{example}

\begin{remark}[Pièges courants du calcul intégral]
(i) \emph{Les \omterm{thm:b1:integration:parts}{changements de variable} doivent être $C^1$ sur tout
l'\omterm{prop:b1:reals:intervals}{intervalle}}~: le changement $x = \frac1t$ est illégal à la
traversée de $0$~; appliqué aveuglément à
$\int_{-1}^{1}\frac{\dd x}{1 + x^2}$, il «~démontre~» que
l'\omterm{thm:b1:integration:def}{intégrale} est égale à son opposée. Lorsqu'un \omterm{thm:b1:integration:parts}{changement de
variable} présente une singularité, on coupe d'abord l'\omterm{prop:b1:reals:intervals}{intervalle}
(Chasles), on substitue sur chaque morceau, et seulement ensuite on
recolle. (ii) \emph{Les primitives logarithmiques réclament des
valeurs absolues}~: $\int \frac{\dd x}{x - 2} = \ln\abs{x - 2} + C$
de chaque côté de $2$ séparément --- écrire $\ln(x - 2)$ sur
$\intoo{0}{1}$, c'est écrire le logarithme d'un nombre négatif~;
et la constante $C$ peut différer de part et d'autre de la
singularité. (iii) \emph{Une \omterm{thm:b1:integration:def}{intégrale} nulle ne tue pas la
fonction}~: $\int_0^{2\pi}\sin = 0$~; la positivité de
l'intégrande est requise avant de conclure $f = 0$
(\cref{ex:b1:integration:chasleswork}). (iv) \emph{Les \omterm{thm:b1:integration:riemann}{sommes de
Riemann} doivent être calibrées}~: dans $\frac{b-a}{n}\sum f\bigl(a +
k\frac{b-a}{n}\bigr)$, le pas à l'extérieur et les points à
l'\omterm{def:b1:topology:closure}{intérieur} doivent correspondre à la même subdivision --- l'erreur
fréquente est une somme $\sum_{k=1}^{n} f\bigl(\frac kn\bigr)$
\emph{sans} le facteur $\frac1n$, qui diverge au lieu de converger
vers $\int_0^1 f$. Liste de contrôle avant d'invoquer
le \cref{thm:b1:integration:riemann}~: mettre $\frac1n$ en facteur,
réécrire le terme général comme $f$ de $\frac kn$, nommer $f$ et
vérifier sa \omterm{def:b1:continuity:continuous}{continuité}.
\end{remark}

\begin{example}[Deviner, dériver, ajuster]\label{ex:b1:integration:guessdiff}
Que vaut $\int_1^x (\ln t)^2\,\dd t$~? Devinons une primitive de
la forme $t\,P(\ln t)$ avec $P$ polynomiale et dérivons~:
\[
\bigl(t\,P(\ln t)\bigr)' = P(\ln t) + P'(\ln t) .
\]
Il faut $P(u) + P'(u) = u^2$~: prenons $P(u) = u^2 - 2u + 2$ (on
identifie les coefficients en descendant depuis $u^2$). D'où
\[
\int_1^x (\ln t)^2\,\dd t
= \bigl[t\bigl((\ln t)^2 - 2\ln t + 2\bigr)\bigr]_1^x
= x(\ln x)^2 - 2x\ln x + 2x - 2 ,
\]
résultat qu'on obtiendrait sinon par deux \omterm{thm:b1:integration:parts}{intégrations par
parties}. L'idée à retenir~: pour des intégrandes de la forme
(\omterm{def:b1:poly:def}{polynôme} en $\ln t$) ou (\omterm{def:b1:poly:def}{polynôme} fois $\eu^{\lambda t}$), la
primitive a la même forme --- dériver une devinette bien formée
convertit l'intégration en algèbre linéaire sur les coefficients,
plus rapide et moins sujette à erreur que des parties itérées.
\end{example}

\begin{example}[L'intégrale boomerang]\label{ex:b1:integration:boomerang}
Calculons $I = \int_0^{\pi/2} \eu^x \cos x\,\dd x$. Intégrons deux
fois par parties, en dérivant à chaque fois le facteur
trigonométrique~:
\[
I = \bigl[\eu^x \sin x\bigr]_0^{\pi/2}
- \int_0^{\pi/2} \eu^x \sin x\,\dd x
= \eu^{\pi/2} - J,
\qquad
J = \bigl[-\eu^x\cos x\bigr]_0^{\pi/2}
+ \int_0^{\pi/2} \eu^x\cos x\,\dd x = 1 + I .
\]
L'\omterm{thm:b1:integration:def}{intégrale} est revenue sur elle-même~: $I = \eu^{\pi/2} - 1 - I$,
d'où
\[
I = \frac{\eu^{\pi/2} - 1}{2} .
\]
L'idée à retenir~: lorsque l'intégrande est le produit de deux
fonctions qui se reproduisent par dérivation ($\eu^{ax}$, $\cos
bx$, $\sin bx$), deux \omterm{thm:b1:integration:parts}{intégrations par parties} produisent une
équation linéaire d'inconnue l'\omterm{thm:b1:integration:def}{intégrale} --- il faut la résoudre
au lieu d'intégrer~; de façon équivalente, on passe par
$\eu^{(a + \iu b)x}$ (\cref{ch:b1:complex}) et on prend la partie
réelle. Les deux chemins donnent la même réponse, et vérifier
qu'ils le font est une vérification gratuite.
\end{example}

\section{Sommes de Riemann}

\begin{theorem}[Sommes de Riemann]\label{thm:b1:integration:riemann}
Soit $f$ \omterm{def:b1:continuity:continuous}{continue} sur $\intcc{a}{b}$. Alors\index{sommes de Riemann}
\[
S_n = \frac{b - a}{n} \sum_{k=0}^{n-1}
f\Bigl(a + k\,\frac{b-a}{n}\Bigr)
\xrightarrow[n \to \infty]{} \int_a^b f(t)\, \dd t ,
\]
et de même avec n'importe quels points d'évaluation à l'intérieur des
\omterm{prop:b1:reals:intervals}{sous-intervalles}.
\end{theorem}

\begin{proof}
$S_n$ est l'\omterm{thm:b1:integration:def}{intégrale} de la \omterm{def:b1:integration:step}{fonction en escalier} $\varphi_n$ égale à
$f(a + k\frac{b-a}{n})$ sur le $k$-ième \omterm{prop:b1:reals:intervals}{sous-intervalle}. Soit
$\varepsilon > 0$~; la \omterm{def:b1:continuity:uniform}{continuité uniforme} (Heine) fournit $\delta$~;
pour $n > \frac{b-a}{\delta}$, tout point d'un \omterm{prop:b1:reals:intervals}{sous-intervalle} est à
distance $< \delta$ de son point d'évaluation, donc $\abs{f -
\varphi_n} \leq \varepsilon$ sur $\intcc{a}{b}$, d'où
\[
\Bigl| \int_a^b f - S_n \Bigr|
= \Bigl| \int_a^b (f - \varphi_n) \Bigr|
\leq (b-a)\,\varepsilon . \qedhere
\]
\end{proof}

\begin{omfigure}
\begin{tikzpicture}
\begin{axis}[omaxis, width=8.6cm, height=5.4cm,
    xmin=0, xmax=1.15, ymin=0, ymax=1.7,
    xtick={0,1}, ytick=\empty]
  \draw[fill=omDef!15, draw=omDef!60, thin]
    (axis cs:0,0)      rectangle (axis cs:0.125,0.4)
    (axis cs:0.125,0)  rectangle (axis cs:0.25,0.4156)
    (axis cs:0.25,0)   rectangle (axis cs:0.375,0.4625)
    (axis cs:0.375,0)  rectangle (axis cs:0.5,0.5406)
    (axis cs:0.5,0)    rectangle (axis cs:0.625,0.65)
    (axis cs:0.625,0)  rectangle (axis cs:0.75,0.7906)
    (axis cs:0.75,0)   rectangle (axis cs:0.875,0.9625)
    (axis cs:0.875,0)  rectangle (axis cs:1,1.1656);
  \addplot[omThm, very thick, domain=0:1.05, samples=60]
    {0.4 + x^2};
\end{axis}
\end{tikzpicture}

{\small Une somme de Riemann à gauche avec $n = 8$ rectangles~:
quand le pas diminue, la \omterm{def:b1:continuity:uniform}{continuité uniforme} force l'aire de
l'escalier vers $\int_a^b f$.}
\end{omfigure}

\begin{example}\label{ex:b1:integration:riemannexample}
$\displaystyle\sum_{k=1}^{n} \frac{1}{n + k} = \frac 1n
\sum_{k=1}^{n} \frac{1}{1 + k/n} \xrightarrow[n\to\infty]{}
\int_0^1 \frac{\dd x}{1 + x} = \ln 2$~: une limite invisible aux
majorations élémentaires, transparente comme somme de Riemann.
\end{example}

\begin{example}[Une deuxième somme de Riemann, avec calibrage]\label{ex:b1:integration:riemann2}
Cherchons
\[
\lim_{n\to\infty} \sum_{k=1}^{n} \dfrac{n}{(n+k)^2} .
\]
Calibrons~:
\[
\sum_{k=1}^{n} \frac{n}{(n+k)^2}
= \frac{1}{n}\sum_{k=1}^{n} \frac{n^2}{(n+k)^2}
= \frac1n \sum_{k=1}^{n} \frac{1}{\bigl(1 + \frac kn\bigr)^2} ,
\]
somme de Riemann de la fonction \omterm{def:b1:continuity:continuous}{continue} $f(x) = \frac{1}{(1+x)^2}$
sur $\intcc{0}{1}$~: la limite est
\[
\int_0^1 \frac{\dd x}{(1 + x)^2}
= \Bigl[-\frac{1}{1+x}\Bigr]_0^1 = \frac12 .
\]
L'idée à retenir~: tout l'art tient dans la ligne du milieu ---
forcer le terme général à prendre la forme $f(\frac kn)$ au prix
de l'extraction d'exactement un facteur $\frac1n$~; une fois la
forme correcte, le théorème fait l'analyse et le théorème
fondamental fait l'arithmétique.
\end{example}

\begin{remark}[Où l'intégrale travaille ensuite]
Chacune des constructions de ce chapitre a une suite. Les \omterm{thm:b1:integration:riemann}{sommes de
Riemann} reviennent au \cref{ch:b1:series} comme pont entre séries
et \omterm{thm:b1:integration:def}{intégrales} (comparaison de $\sum \frac{1}{n^\alpha}$ avec $\int
\frac{\dd t}{t^\alpha}$)~; le reste intégral est la forme la plus
fine de la formule de Taylor au \cref{ch:b1:taylor}~; la définition
par $\sup$ est le prototype de l'\omterm{thm:b1:integration:def}{intégrale} de Lebesgue du volume de
Licence 3, où les trois mêmes propriétés (linéarité, croissance, un
théorème de convergence) sont rebâties sur une classe de fonctions
bien plus large. Et le devoir maison ci-dessous transforme
l'\omterm{thm:b1:integration:parts}{intégration par parties} en arithmétique~: l'irrationalité de
$\pi^2$.
\end{remark}

\section{Exercices}

\begin{exercise}[$\star$]\label{exo:b1:integration:1}
Calculer~: $\displaystyle\int_0^1 \frac{\dd x}{x^2 - 4}$
\emph{(éléments simples, \cref{ch:b1:fractions})}~;
$\displaystyle\int_1^{\eu} \ln t \,\dd t$~;
$\displaystyle\int_0^{\pi} t \sin t\, \dd t$.
\end{exercise}

\begin{exercise}[$\star$]\label{exo:b1:integration:2}
Calculer $\displaystyle\int_0^{1} \frac{t}{(t^2+1)^2}\,\dd t$
(\omterm{thm:b1:integration:parts}{changement de variable}) et
$\displaystyle\int_0^{\pi/2} \cos^3 t\, \dd t$ (écrire $\cos^3 =
\cos(1 - \sin^2)$).
\end{exercise}

\begin{exercise}[$\star$]\label{exo:b1:integration:3}
Déterminer les limites, comme \omterm{thm:b1:integration:riemann}{sommes de Riemann}~:
\[
u_n = \sum_{k=1}^{n} \frac{n}{n^2 + k^2},
\qquad
v_n = \frac{1}{n}\sqrt[n]{\frac{(2n)!}{n!\,n^n}}
\ \emph{(prendre le logarithme)}.
\]
\end{exercise}

\begin{exercise}[$\star$]\label{exo:b1:integration:4}
Soit $f$ \omterm{def:b1:continuity:continuous}{continue} sur $\intcc{0}{1}$. Calculer $\lim_{n\to\infty}
\int_0^1 x^n f(x)\,\dd x$. \emph{(Couper $\intcc{0}{1}$ en $1 -
\delta$.)}
\end{exercise}

\begin{exercise}[$\star\star$]\label{exo:b1:integration:5}
(Cauchy--Schwarz) Pour $f, g$ \omterm{def:b1:continuity:continuous}{continues} sur $\intcc{a}{b}$, démontrer
\[
\Bigl(\int_a^b fg\Bigr)^{\!2} \leq \int_a^b f^2 \cdot \int_a^b g^2 ,
\]
en développant $\int_a^b (f + \lambda g)^2 \geq 0$ comme un trinôme en
$\lambda$. Quand y a-t-il égalité~?
\end{exercise}

\begin{exercise}[$\star\star$]\label{exo:b1:integration:6}
Soit $f$ \omterm{def:b1:continuity:continuous}{continue} sur $\R$, $T$-périodique. Démontrer que
$\int_a^{a+T} f$ ne dépend pas de $a$, et que $\frac1x \int_0^x
f(t)\,\dd t \to \frac 1T \int_0^T f$ quand $x \to +\infty$.
\end{exercise}

\begin{exercise}[$\star\star$]\label{exo:b1:integration:7}
Pour $f$ \omterm{def:b1:continuity:continuous}{continue} sur $\intcc{0}{1}$ avec $\int_0^1 f = \frac12$,
démontrer que $f$ admet un point fixe dans $\intcc{0}{1}$.
\emph{(Intégrer $f(x) - x$ et utiliser la positivité stricte,
\cref{thm:b1:integration:props}~(4), par sa contraposée combinée au
\omterm{thm:b1:continuity:ivt}{théorème des valeurs intermédiaires}.)}
\end{exercise}

\begin{exercise}[$\star\star\star$]\label{exo:b1:integration:8}
(\omterm{thm:b1:integration:def}{Intégrales} de Wallis) Soit $W_n = \int_0^{\pi/2} \sin^n t\,\dd t$.
\begin{enumerate}
  \item Démontrer la relation de récurrence $n W_n = (n-1) W_{n-2}$
        ($n \geq 2$) par parties, et calculer $W_0, W_1$, puis
        $W_{2p}$ et $W_{2p+1}$ sous forme close.
  \item Démontrer que $(W_n)$ est décroissante avec
        $\frac{W_{n+1}}{W_n} \to 1$~; démontrer que la quantité
        $(n+1)\,W_{n+1} W_n$ est constante, égale à $\frac\pi2$~; et
        en déduire l'équivalent $W_n \sim \sqrt{\dfrac{\pi}{2n}}$.
\end{enumerate}
\end{exercise}

\begin{exercise}[$\star\star\star$]\label{exo:b1:integration:9}
(Niven~: $\pi$ est irrationnel) Supposons $\pi = \frac ab$ avec $a, b
\in \N^*$, et posons, pour $n$ à choisir,
\[
P(x) = \frac{x^n (a - bx)^n}{n!},
\qquad
I_n = \int_0^{\pi} P(x)\sin x\, \dd x .
\]
\begin{enumerate}
  \item Démontrer que $0 < I_n \leq \pi\,\frac{(\pi a)^n}{n!}$, qui
        est $< 1$ pour $n$ grand.
  \item Démontrer que $P$ et toutes ses \omterm{def:b1:derivative:def}{dérivées} prennent des valeurs
        \emph{entières} en $0$ et en $\pi = \frac ab$.
        \emph{(Développement du binôme~: les coefficients de $P$
        multipliés par $k!$ sont entiers pour $k \geq n$~; et
        $P(\pi - x) = P(x)$.)}
  \item Poser $Q = P - P'' + P^{(4)} - \dots$ (somme finie). Vérifier
        que $\bigl(Q'\sin x - Q\cos x\bigr)' = P \sin x$, et en
        déduire que $I_n = Q(\pi) + Q(0)$ est un entier.
  \item Conclure.
\end{enumerate}
\end{exercise}

\begin{exercise}[$\star\star\star$]\label{exo:b1:integration:10}
Soit $f$ de classe $C^1$ sur $\intcc{a}{b}$. Démontrer la limite de
type Riemann--Lebesgue
\[
\int_a^b f(t)\sin(\lambda t)\,\dd t \xrightarrow[\lambda \to
+\infty]{} 0
\]
en intégrant par parties. Puis la démontrer à nouveau pour $f$
seulement \omterm{def:b1:continuity:continuous}{continue}, par approximation uniforme par des \omterm{def:b1:integration:step}{fonctions en
escalier} (\cref{thm:b1:integration:approx}).
\end{exercise}

\begin{exercise}[$\star\star$]\label{exo:b1:integration:11}
(Formule de la moyenne) Soient $f, g$ \omterm{def:b1:continuity:continuous}{continues} sur $\intcc{a}{b}$
avec $g \geq 0$. Démontrer qu'il existe $c \in \intcc{a}{b}$ tel que
\[
\int_a^b f(t)\,g(t)\,\dd t = f(c)\int_a^b g(t)\,\dd t ,
\]
et montrer par un exemple que l'hypothèse $g \geq 0$ ne peut pas être
supprimée.
\end{exercise}

\begin{exercise}[$\star\star\star$]\label{exo:b1:integration:12}
(Les moments forcent des zéros) Soit $f$ \omterm{def:b1:continuity:continuous}{continue} sur $\intcc{a}{b}$
avec
\[
\int_a^b f(t)\,t^k\,\dd t = 0 \qquad \text{pour } k = 0, 1,
\dots, n .
\]
Démontrer que $f$ s'annule en $n + 1$ points distincts de
$\intoo{a}{b}$. \emph{(Si $f$ ne change de signe qu'en $z_1 < \dots
< z_m$ avec $m \leq n$, intégrer $f$ contre $P(t) = (t - z_1)
\cdots (t - z_m)$ et utiliser la positivité stricte.)}
\end{exercise}

\begin{remark}[Perspectives à l'intérieur de ce volume]
Trois chapitres à venir s'appuient directement sur celui-ci. Le
\cref{ch:b1:taylor} porte le reste intégral --- la plus fine des
trois formules de Taylor est une \omterm{thm:b1:integration:parts}{intégration par parties} itérée $n$
fois. Le \cref{ch:b1:series} convertit l'encadrement des sommes par
des \omterm{thm:b1:integration:def}{intégrales} en le critère décisif pour $\sum n^{-\alpha}$, et son
devoir maison affine cet encadrement en la constante d'Euler. Le
\cref{ch:b1:curves} rend l'\omterm{thm:b1:integration:def}{intégrale} géométrique~: la longueur d'un
arc paramétré est $\int \sqrt{x'(t)^2 + y'(t)^2}\,\dd t$, \omterm{thm:b1:integration:def}{intégrale}
d'une fonction \omterm{def:b1:continuity:continuous}{continue} sur un segment --- précisément l'objet
construit ici, sans aucune théorie des \omterm{thm:b1:integration:def}{intégrales} impropres. Le fait
le plus réutilisé sera le plus humble~: $\bigl|\int f\bigr| \leq (b
- a)\sup\abs{f}$, l'inégalité qui transforme toute estimation
ponctuelle en estimation \omterm{thm:b1:integration:def}{intégrale}.
\end{remark}

\section{Problème~: la machine à irrationalité par les intégrales}

\begin{problem}[{Devoir maison --- $\eu$ et $\pi^2$ sont
irrationnels, $\eu$ à six décimales, et $\frac{22}{7} > \pi$ avec
démonstration}]
\label{pb:b1:integration:1}
Un seul mécanisme alimente tout ce problème~: \emph{une expression
qui doit être un entier strictement positif, tout en étant
démontrablement plus petite que $1$, ne peut pas exister}.
L'\cref{exo:b1:integration:9} (Niven) l'a fait tourner une fois pour
démontrer $\pi \notin \Q$~; ici, nous l'industrialisons. La machine
a besoin de trois pièces~: une entrée d'\emph{intégralité} (les
valeurs aux bornes de \omterm{def:b1:poly:def}{polynômes} bien choisis), une entrée de
\emph{petitesse} (un facteur $\frac{1}{n!}$ qui écrase l'\omterm{thm:b1:integration:def}{intégrale}),
et un \emph{pont} (l'\omterm{thm:b1:integration:parts}{intégration par parties}) qui les relie. Nous
démontrons que $\eu$ est irrationnel et le calculons avec une erreur
certifiée, démontrons le théorème plus fin de Legendre selon lequel
$\pi^2$ est irrationnel, et terminons par la plus charmante
\omterm{thm:b1:integration:def}{intégrale} de l'analyse~: $\int_0^1 \frac{x^4(1 - x)^4}{1 + x^2}\dd x
= \frac{22}{7} - \pi$, qui encadre $\pi$ à la
main.\index{irrationalité!de pi carré@de
$\pi^2$}\index{irrationalité!de e@de $\eu$}

\medskip
\noindent\textbf{Partie I --- le carburant.}
\begin{enumerate}
  \item Démontrer que $\dfrac{c^{\,n}}{n!} \to 0$ pour tout $c > 0$
        fixé \emph{(au-delà de $n \geq 2c$, chaque étape \omterm{def:b1:arith:divides}{divise} le
        terme au moins par deux)}.
  \item (\omterm{thm:b1:integration:def}{Intégrales} bêta) Démontrer, par récurrence sur $m$ avec une
        \omterm{thm:b1:integration:parts}{intégration par parties}~:
        \[
        \int_0^1 x^{\,k}\,(1 - x)^{\,m}\,\dd x =
        \frac{k!\,m!}{(k + m + 1)!}
        \qquad (k, m \in \N).
        \]
  \item En déduire $\displaystyle\int_0^1 \bigl(x(1-x)\bigr)^n \dd x
        = \frac{1}{(2n+1)\binom{2n}{n}}$, et --- en comparant avec
        la majoration $x(1 - x) \leq \frac14$ --- l'estimation
        $\binom{2n}{n} \geq \dfrac{4^n}{2n+1}$, en accord avec
        $\binom{2n}{n}^{1/n} \to 4$ du \cref{pb:b1:seq:1}.
  \item Démontrer le lemme de petitesse utilisé deux fois
        ci-dessous~: pour toute fonction \omterm{def:b1:continuity:continuous}{continue} $g > 0$ sur
        $\intoo{0}{1}$,
        \[
        0 < \int_0^1 \bigl(x(1-x)\bigr)^n g(x)\,\dd x
        \leq \frac{\sup_{\intcc{0}{1}}\abs g}{4^{\,n}} .
        \]
\end{enumerate}

\medskip
\noindent\textbf{Partie II --- $\eu$~: irrationalité, puis six
décimales.} Posons $A_n = \displaystyle\int_0^1 x^n \eu^x \dd x$.
\begin{enumerate}[resume]
  \item Calculer $A_0$ et $A_1$, démontrer la relation de récurrence
        $A_n = \eu - n A_{n-1}$, et l'encadrement $0 < A_n \leq
        \dfrac{\eu}{n+1}$.
  \item Montrer par récurrence que $A_n = \alpha_n + \beta_n \eu$
        avec $\alpha_n, \beta_n \in \Z$.
  \item En déduire que $\eu$ est irrationnel \emph{(si $\eu = \frac
        pq$, alors $q A_n$ est un entier piégé dans
        $\intoo{0}{1}$ pour $n$ grand)}. Comparer avec la
        démonstration de l'\cref{exo:b1:seq:9}~: même chute,
        carburant différent.
  \item Démontrer, par récurrence et \omterm{thm:b1:integration:parts}{intégration par parties}, la
        formule exacte du reste
        \[
        \eu = \sum_{k=0}^{n} \frac{1}{k!} + R_n,
        \qquad
        R_n = \frac{1}{n!}\int_0^1 (1 - t)^n\,\eu^{\,t}\,\dd t,
        \qquad
        \frac{1}{(n+1)!} \leq R_n \leq \frac{\eu}{(n+1)!} .
        \]
  \item Prendre $n = 9$~: majorer $R_9$ à l'aide de $\eu < 2.75$
        (issu de $b_2 = 2.75$ dans l'\cref{ex:b1:seq:e}), évaluer la
        somme, et conclure l'encadrement certifié $2.7182818 \leq
        \eu \leq 2.7182823$ --- six décimales, $\eu \approx
        2.718282$, avec démonstration.
\end{enumerate}

\medskip
\noindent\textbf{Partie III --- le théorème de Legendre~: $\pi^2$ est
irrationnel.} Soit $f(x) = \dfrac{x^n (1 - x)^n}{n!}$, et supposons
$\pi^2 = \frac ab$ avec $a, b \in \N^*$.
\begin{enumerate}[resume]
  \item Montrer que $f(1 - x) = f(x)$ et $0 < f \leq
        \dfrac{1}{4^n\,n!}$ sur $\intoo{0}{1}$.
  \item Montrer que $f^{(k)}(0)$ et $f^{(k)}(1)$ sont des entiers
        pour tout $k \geq 0$ \emph{(développer $x^n(1-x)^n$ à
        coefficients entiers~; $\frac{k!}{n!} \in \Z$ pour $k \geq
        n$~; puis utiliser la symétrie)}.
  \item Définir
        \[
        G = b^{\,n} \sum_{k=0}^{n} (-1)^k\,
        \pi^{2n - 2k} f^{(2k)} .
        \]
        Montrer que $G(0)$ et $G(1)$ sont des entiers \emph{(chaque
        $b^n \pi^{2n-2k} = a^{\,n-k}\,b^{\,k}$)}.
  \item Vérifier le télescopage $G'' + \pi^2 G = b^n \pi^{2n+2} f
        = \pi^2 a^n f$, puis
        \[
        \frac{\dd}{\dd x}\Bigl(G'(x)\sin \pi x - \pi\,G(x)\cos
        \pi x\Bigr) = \pi^2 a^n f(x)\sin \pi x .
        \]
  \item Intégrer sur $\intcc{0}{1}$ et conclure
        \[
        \pi\,a^n \int_0^1 f(x)\sin(\pi x)\,\dd x = G(0) + G(1)
        \in \Z ,
        \]
        un entier \emph{strictement positif} majoré par $\dfrac{\pi
        a^n}{4^n\,n!}$.
  \item Conclure avec la question 1 que $\pi^2$ est irrationnel
        (Legendre, 1794), et que cela renforce
        l'\cref{exo:b1:integration:9}~: pourquoi l'irrationalité de
        $\pi^2$ implique-t-elle celle de $\pi$, et non l'inverse~?
\end{enumerate}

\medskip
\noindent\textbf{Partie IV --- comprendre la machine.}
\begin{enumerate}[resume]
  \item Repérer les deux forces opposées (l'intégralité des données
        aux bornes~; la petitesse analytique de l'\omterm{thm:b1:integration:def}{intégrale}) et le
        pont, dans les parties II et III. Expliquer ensuite pourquoi
        le facteur $\frac{1}{n!}$ dans $f$ est le nœud de
        l'affaire~: si on le retire, l'intégralité subsiste, mais
        quelle inégalité meurt, et pour quelles fractions $\frac ab$
        prétendues la démonstration échoue-t-elle alors~?
  \item Effectivité~: supposons que quelqu'un prétende $\pi^2 =
        \frac ab$ avec $a \leq 10$. Montrer que la contradiction
        tombe déjà à $n = 7$~: calculer $\pi\,(10/4)^7/7! \approx
        0.38 < 1$. La machine ne se contente pas de réfuter~; elle
        réfute à un rang fixé et calculable.
  \item Vérifier le pont sans aucune hypothèse~: démontrer par deux
        \omterm{thm:b1:integration:parts}{intégrations par parties} que
        \[
        \int_0^1 x(1 - x)\sin(\pi x)\,\dd x = \frac{4}{\pi^3},
        \]
        et recoller avec la question 14 pour $n = 1$ (garder
        $\pi^2$ symbolique~: l'identité télescopée s'écrit $\pi^3
        \int_0^1 f_1 \sin \pi x = -(f_1''(0) + f_1''(1)) = 4$).
  \item Qu'est-ce qui rend $\eu^x$ et $\sin \pi x$ éligibles comme
        noyaux de la machine~? Identifier la propriété (chacun
        vérifie une équation différentielle linéaire à coefficients
        constants, de sorte que les \omterm{thm:b1:integration:parts}{intégrations par parties}
        répétées bouclent sur le point de départ), et nommer la
        \omterm{def:b1:topology:closure}{frontière}~: la même machine, affinée par Hermite et
        Lindemann, démontre que $\eu$ et $\pi$ sont
        \emph{\omterm{pb:b1:logic:1}{transcendants}} --- au-delà de ce volume.
\end{enumerate}

\medskip
\noindent\textbf{Partie V --- $\frac{22}{7}$ contre $\pi$, et la
morale.}
\begin{enumerate}[resume]
  \item Établir la division polynomiale
        \[
        \frac{x^4(1-x)^4}{1 + x^2}
        = x^6 - 4x^5 + 5x^4 - 4x^2 + 4 - \frac{4}{1 + x^2},
        \]
        et en déduire la célèbre identité
        \[
        \int_0^1 \frac{x^4 (1-x)^4}{1 + x^2}\,\dd x
        = \frac{22}{7} - \pi .
        \]
  \item L'intégrande est positive~: en conclure $\pi <
        \frac{22}{7}$. Puis, en encadrant $\frac{1}{1+x^2}$ entre
        $\frac12$ et $1$ et en utilisant $\int_0^1 (x(1-x))^4 =
        \frac{1}{630}$ (question 3), démontrer
        \[
        \frac{22}{7} - \frac{1}{630} \;\leq\; \pi \;\leq\;
        \frac{22}{7} - \frac{1}{1260} ,
        \]
        c'est-à-dire $3.14126 \leq \pi \leq 3.14207$~: deux
        décimales exactes, à la main.
  \item Généraliser~: en divisant $x^{4m}(1-x)^{4m}$ par $1 + x^2$,
        montrer que le reste est la constante $(-4)^m$
        \emph{(travailler modulo $x^2 + 1$~: $(1-x)^2 \equiv -2x$)},
        en déduire des rationnels $r_m$ tels que
        \[
        \abs{\pi - r_m} \leq 4^{\,1 - 5m} ,
        \]
        et vérifier que $m = 1$ redonne les questions 20--21.
  \item Confronter ces rationnels à la théorie de l'approximation du
        \cref{pb:b1:derivative:1}~: calculer $\abs{\pi -
        \frac{22}{7}} \approx 1.26\cdot10^{-3}$ face à la garantie
        de Dirichlet $\frac{1}{49}$, et citer
        $\abs{\pi - \frac{355}{113}} \approx 2.7\cdot10^{-7}$
        face à $\frac{1}{113^2} \approx 7.8\cdot10^{-5}$~: il
        existe des approximations rationnelles exceptionnellement
        bonnes de $\pi$ --- ce qui est cohérent, puisqu'on ne sait
        pas que $\pi$ soit mal approchable.
  \item (Le piège de l'entier, en abstrait) Démontrer le lemme qui
        unifie tout~: \emph{si $x \in \R$ et s'il existe des entiers
        $a_n, b_n$ tels que $0 < \abs{a_n + b_n x} \to 0$, alors $x$
        est irrationnel.} Recenser ses occurrences dans ce problème,
        dans l'\cref{exo:b1:integration:9}, dans
        l'\cref{exo:b1:seq:9}, et dans le \cref{pb:b1:derivative:1}.
  \item Synthèse, une phrase pour chacun~: (i) les trois pièces de
        la machine et l'endroit où chacune vit dans la boîte à
        outils de ce chapitre~; (ii) ce que l'\omterm{thm:b1:integration:def}{intégrale} apporte que
        le \omterm{thm:b1:derivative:mvt}{théorème des accroissements finis} du
        \cref{pb:b1:derivative:1} ne pouvait pas apporter~; (iii)
        l'inventaire des résultats extraits (deux irrationalités,
        une constante à six décimales, un encadrement de $\pi$, une
        minoration binomiale)~; (iv) la \omterm{def:b1:topology:closure}{frontière} (Hermite,
        Lindemann~; et le même piège, appliqué à $\zeta(2)$ et
        $\zeta(3)$, dans l'arithmétique du vingtième siècle).
\end{enumerate}
\end{problem}
